\documentclass[a4paper, 12pt]{article}
\usepackage[utf8]{inputenc}
\usepackage[spanish, provide=*]{babel}
\usepackage{amsmath, amssymb}
\usepackage{graphicx}
\usepackage{hyperref}
\usepackage{listings}
\usepackage{tcolorbox}
\usepackage{xcolor}
\usepackage{fancyhdr}
\usepackage{geometry}

\geometry{left=2.5cm, right=2.5cm, top=3cm, bottom=4.5cm}

% Para las tablas
\usepackage{float}  %Para que no se cambien de lugar
\usepackage{pgfplotstable}
\usepackage{array}
\usepackage{colortbl}
\usepackage{url} % para que el campo note muestre URLs bien

% Estilo
\usepackage{lastpage}

% Definición de colores
\definecolor{cafetitulos}{RGB}{150,75,0}
\definecolor{redv}{RGB}{165, 43, 42}
\definecolor{cafe}{RGB}{130,75,0}
\definecolor{cafet}{RGB}{204, 121, 5}

% Configuración de encabezados y pies de página
\pagestyle{fancy}
\renewcommand{\headrulewidth}{0pt} % Elimina la línea por defecto del encabezado

% Encabezado
\fancyhead[L]{\includegraphics[height=1.5cm]{logomanual.png}}
\fancyhead[R]{\nouppercase{\leftmark}} % Nombre de la sección

% Pie de página
\renewcommand{\footrulewidth}{0.5pt} % Línea separadora
\renewcommand{\footrule}{\hbox to\headwidth{\color{cafe}\leaders\hrule height \footrulewidth\hfill}}
\fancyfoot[L]{Manual Técnico - 2025}
\fancyfoot[C]{}
\fancyfoot[R]{Página \thepage} % Número de página a la derecha

% Ajustar altura del encabezado y espacio para el pie
\setlength{\headheight}{52pt} % Altura del encabezado (ajusta según el logo)
\setlength{\footskip}{40pt} % Espacio entre el texto y el pie (evita solapamiento)

% Titulos
\usepackage{titlesec} % Para personalizar section/subsection
\titleformat{\section}
  {\normalfont\Large\bfseries}{\thesection}{1em}{\color{cafetitulos}} % Section en café
\titleformat{\subsection}
  {\normalfont\bfseries}{\thesubsection}{1em}{\color{cafe}}
\titleformat{\subsubsection}{\normalfont\bfseries}{\thesubsubsection}{1em}{\color{cafetitulos}}

% Configuración para código fuente
\lstset{
    language=Java, % puedes cambiar a Python, SQL, etc.
    basicstyle=\ttfamily\small,
    backgroundcolor=\color{gray!10},
    frame=single,
    breaklines=true,
    keywordstyle=\color{cafetitulos},
    commentstyle=\color{gray},
    stringstyle=\color{redv}
}


\title{
\includegraphics[width=5cm]{logo.png}\\

{\Large CHIBCHAWEB HOSTING PLATFORM}\\
{\Large APLICACIÓN PARA UNA COMPAÑÍA DE HOSPEDAJE DE PÁGINAS WEB}\\[0.5em]

\textbf{MANUAL TÉCNICO}\\[0.5em]
}
\author{
\textbf{Grupo A}\\[0.5em]
Ingenieros:\\
Andrés José Acevedo Ardila\\
Brayan Steven Sánchez Casilima\\
Christian Camilo Lancheros Sánchez\\
Juan Esteban Ordoñez Sandoval\\
Karen Tatiana Bravo Rodríguez\\
Luis Felipe Mayorga Tibaquicha\\\\\\
Universidad Distrital Francisco José de Caldas\\
Facultad de Ingeniería\\
Ingeniería de sistemas\\
}
\date{Bogotá D.C., Agosto 2025}


% Inicio Documento
\begin{document}
\maketitle
\newpage

\begin{titlepage}
\newgeometry{left=2cm, right=0cm, top=2cm, bottom=2cm} % sin margen derecho
\thispagestyle{empty}

\begin{tikzpicture}[remember picture, overlay]
    \node[anchor=north east, inner sep=0pt, xshift=12cm] at (current page.north east) {
        \includegraphics[height=\paperheight]{registro-imagen.png}
    };
\end{tikzpicture}

\vfill
{\includegraphics[width=0.6\textwidth]{logomanual.png}\par}
\vfill
{\Huge Grupo A\par}
{\Huge Manual Técnico ChibchaWeb\par}
{\Large 2025\par}
\vfill


\end{titlepage}


\newpage
\tableofcontents

\newpage
\listoffigures

\newpage
\listoftables

\newpage
\section{Introducción}
Este manual proporciona la información necesaria para la instalación, configuración, operación y mantenimiento del sistema ChibchaWeb, una plataforma web orientada al servicio de hospedaje de sitios web. El sistema permite la gestión de cuentas de clientes, distribuidores y empleados, la adquisición de dominios, facturación automática, administración de tickets de soporte, y más. Esta guía está dirigida a desarrolladores, administradores del sistema y personal técnico encargado del mantenimiento de la solución. \\

El manual técnico documenta el desarrollo e instalación de un sistema web que automatiza las operaciones clave de ChibchaWeb, incluyendo: 

\begin{enumerate}
    \item {Registro y gestión de cuentas de clientes, empleados y distribuidores.}
    \item {Validación de información sensible como direcciones y tarjetas de crédito.}
    \item {Procesamiento de pagos mediante tarjetas aceptadas (VISA, MASTERCARD, DINERS).}
    \item {Administración de dominios y envío de solicitudes a registradores externos.}
    \item {Registro y gestión de tickets de soporte técnico.}
    \item {Cálculo de comisiones para distribuidores y generación de archivos de transferencia. }
\end{enumerate}

El objetivo de esta documentación es proporcionar una guía técnica completa que facilite la instalación, configuración, uso, mantenimiento y actualización del sistema para desarrolladores, administradores del sistema y personal técnico de soporte. 

\newpage
\section{Alcance del manual}
Este documento está dirigido a desarrolladores, administradores de sistemas y personal de soporte que requieran instalar, configurar, mantener y operar la plataforma de gestión de servicios de ChibchaWeb.
El alcance del manual incluye instrucciones detalladas para:

\begin{enumerate}
    \item {La instalación del entorno de desarrollo y producción.}
    \item {La configuración de los archivos clave del sistema.}
    \item {La conexión con la base de datos y los servicios externos.}
    \item {La descripción funcional de cada componente del sistema.}
    \item {Las operaciones necesarias para ejecutar y mantener la plataforma.}
\end{enumerate}


El sistema está compuesto por una arquitectura web basada en tecnologías modernas que permiten escalabilidad, claridad en la gestión de roles y eficiencia en la administración de dominios, pagos y soporte técnico.

\subsection{Frontend}
\begin{itemize}
    \item \textbf{Framework}: Flask (utilizado también como microservidor web para el desarrollo). \cite{flask}
    \item \textbf{Estilos y UI:} Bootstrap. \cite{bootstrap}
\end{itemize}


\subsection{Backend}
\begin{itemize}
    \item \textbf{Lenguaje}: Python.
    \item \textbf{Framework principal}: Flask
\end{itemize}


\subsection{Base de Datos}
\begin{itemize}
    \item \textbf{MySQL} gestionada y visualizada a través de DBeaver como cliente de administración. \cite{mysql}
\end{itemize}


\subsection{Control de Versiones y Colaboración}
\begin{itemize}
	\item \textbf{GitHub:} para manejo de versiones, control de ramas y colaboración entre el equipo.
\end{itemize}

\subsection{Servicios Externos}
\begin{itemize}
	\item \textbf{OpenAI API:} para la automatización en la respuesta a tickets de soporte técnico.
\end{itemize}


\newpage
\section{Requisitos}
\subsection{Requisitos del Hardware}
\begin{itemize}
    \item {\textbf{RAM:} 4GB mínimo, 16 GB ideal.}
    \item { \textbf{CPU}: 2 o más núcleos.}
    \item {\textbf{Almacenamiento}: 4 GB  mínimo, 10 GB ideal.}
\end{itemize}

\subsection{Requisitos de Software}
\begin{itemize}
    \item \textbf{Sistema operativos:}
        \begin{itemize}
            \item {Windows 10/11.}
            \item {Linux}
        \end{itemize}
    \item \textbf{Navegador}: Para que el aplicativo web en desarrollo funcione correctamente el navegador que el usuario use debe tener como requisito las siguientes características:
    \begin{itemize}
        \item {Soporte JavaScript ES6+: React, JSX.}
        \item {Soporte Web APIs modernas.}
        \item {Soporte CSS: Grid/Flexbox.}
        \item {Soporte HTML5.}
    \end{itemize}
    Como recomendación se dará un listado de la versión mínima recomendada para los navegadores más usados:
     \begin{itemize}
        \item {Google Chrome: 100+}
        \item Mozilla: 100+
        \item Microsoft Edge: 100+
        \item Safari (macOS): 15+
        \item Opera: 85+
        \item Safari (iOS): 15+
        \item Chrome/Edge Android: Última versión
    \end{itemize}
    \item \textbf{Entorno de desarrollo}
    \begin{itemize}
        \item Node.js: v18.
        \item npm: v9.
        \item Python: v3.12.
        \item PIP: con las siguientes librerías para en Python.
        \begin{itemize}
            \item PIP
            \item fastapi
            \item Uvicorn
            \item requests
            \item beautifulsoup4
            \item sqlalchemy
            \item pymysql
            \item python-dotenv
            \item email-validator
            \item bcrypt
            \item passlib
            \item cryptography
            \item reportlab
            \item Openai
        \end{itemize}
        \item IDE:
        \begin{itemize}
            \item VS Code
            \item Zed
        \end{itemize}
        \item Base de datos
        \begin{itemize}
            \item MySQL
        \end{itemize}
        \item Gestor base de datos:
        \begin{itemize}
            \item DBeaver
        \end{itemize}
        \item {Bundle}
        \begin{itemize}
            \item Vite
        \end{itemize}
        \item Control de versiones:
        \begin{itemize}
            \item Git
        \end{itemize}
    \end{itemize}
\end{itemize}


\newpage
\section{Instalación}
Para realizar la instalación se recomenda crear una carpeta general para el proyecto, dentro de esta carpetas separadas por cada lógica.

\subsection{Base de datos}
Para realizar la instalación de la base de datos se cuenta con las siguientes dos opciones:
\begin{itemize}
    \item Instalación local la cual permite realizar cambios que no afectan la lógica de desarrollo ni la de producción.
    \item Instalación producción la cual permite actualizar la base de datos que se comunica con todos los clientes del aplicativo.
\end{itemize}

Esta instalación es la que debe usarse para las diferentes etapas del desarollo de cualquier nueva funcionalidad.

Para realizarse se deben seguir los siguientes procesos:

\subsubsection*{Instalación MySQL}
\begin{enumerate}
	\item Instalar MySQL Server.
	\item Crear la base de datos.

	\begin{figure}
	\begin{lstlisting}[language=sql]
CREATE DATABASE chibchaweb;
        \end{lstlisting}
        \caption{Crear base de datos.}
        \label{fig:crear-bd}
        \end{figure}
    \item Crear un nuevo usuario.

    \begin{figure}
    \begin{lstlisting}[language=sql]
CREATE USER 'chibcha'@'localhost' identified by 'chibcha';
        \end{lstlisting}
        \caption{Crear usuario en base de datos.}
        \label{fig:crear-usuario}
        \end{figure}
	\item Otorgar y actualizar privilegios al usuario.

	\begin{figure}
	\begin{lstlisting}[language=sql]
GRANT ALL PRIVILEGES ON chibcha.* TO 'chibcha'@'localhost';
FLUSH PRIVILEGES;
        \end{lstlisting}
        \caption{Dar privilegios al usuario.}
        \label{fig:dar-privilegios}
        \end{figure}
\end{enumerate}

\subsubsection*{Configuración DBeaver}
Ya instalado el Dbeaver se procede a crear nueva conexión en DBeaver:
\begin{enumerate}
\item Abrir DBeaver.
\item Clic en "New Database Connection" (ícono de plug).
\item Eligir MySQL.
\item Llena los campos:
    \begin{itemize}
    \item Host:	localhost
    \item Port:	3306
    \item Database:	chibchaweb
    \item Username:	chibcha
    \item Password:	chibcha
    \end{itemize}
\item Clic en Test Connection.
    \begin{itemize}
	\item Permitir instalar el driver JDBC, solo aceptar y dejar que DBeaver lo descargue.
    \end{itemize}
\item Clic en Finish.
\end{enumerate}

\subsubsection*{Configuración Base de datos}
\begin{enumerate}
	\item Clonar el repositorio:

	\begin{figure}
	\begin{lstlisting}[language=bash]
git clone https://github.com/Rocnarx/Chibchaweb-Hosting-Platform-Database.git
    \end{lstlisting}
    \caption{Clonar repo Database.}
    \label{fig:clonar-db}
    \end{figure}
     \item Ejecutar en Dbeaver el archivo CreacionDB.sql
     \item Ejecutar en DBeaver el archivo DatosBasicos.sql

\end{enumerate}

\subsection{Backend}
\subsubsection*{Configurar requirimientos}
\begin{enumerate}
	\item Descargar e instalar Python 3.12
	\begin{itemize}
	\item Windows:
	Descargar de https://www.python.org/downloads/release/python-3120/
	\item Linux:

	\begin{figure}
	\begin{lstlisting}[language=bash]
sudo apt update
sudo apt install software-properties-common
sudo add-apt-repository ppa:deadsnakes/ppa
sudo apt update
sudo apt install python3.12 python3.12-venv python3.12-dev
        \end{lstlisting}
        \caption{Instalar Python 3.12}
        \label{fig:install-python}
        \end{figure}
	\end{itemize}
\end{enumerate}

\subsubsection*{Configurar Back}
\begin{enumerate}
    \item En la carpeta back crear un entorno virtual

    \begin{figure}
    \begin{lstlisting}[language=bash]
    python3 -m venv ChibchaBack
       \end{lstlisting}
       \caption{Crear venv.}
       \label{fig:crear-venv}
       \end{figure}
    \item Iniciar el entorno virtual


    \begin{figure}
    \begin{lstlisting}[language=bash]
    # -- Windows
    ./ChibchaBack/Scripts/activate
    #Linux
    source ChibchaBack/bin/activate
    \end{lstlisting}
    \caption{Activar venv}
    \label{fig:Activar-venv}
    \end{figure}

	\item Clonar el repositorio:

	\begin{figure}
	\begin{lstlisting}[language=bash]
    git clone https://github.com/Rocnarx/Chibchaweb-Hosting-Platform-Backend.git
    \end{lstlisting}
    \caption{Clonar repo backend.}
    \label{fig:Clonar-repo-backend}
    \end{figure}

    \item Instalar los paquetes requeridos que se encuentran en archivo requeriments.text:


    \begin{figure}
    \begin{lstlisting}[language=bash]
    pip install -r Chibchaweb-Hosting-Platform-Backend/requirements.txt
    \end{lstlisting}
    \caption{Instalar requerimientos back.}
    \label{fig:install-req}
    \end{figure}

    \item Crear archivo .env el cual nos permetira conectar el back con los demás componentes

    \item Configurar archivo .env con el siguiente texto

    \begin{figure}
    \begin{lstlisting}[language=bash]
FERNET_KEY= Solicitar al equipo de desarrollo
API_KEY= Solicitar al equipo de desarrollo
DATABASE_URL= mysql+pymysql://root:.. Solicitar al equipo de desarrollo
OPENAI_API_KEY= Solicitar al equipo de desarrollo
\end{lstlisting}
\caption{Editar .env back}
\label{fig:env-back}
\end{figure}

\end{enumerate}

\subsection{Front}

\subsubsection*{Configurar requerimientos}
\begin{enumerate}
    \item Instalar Node.js y npm
    \begin{itemize}
        \item Windows: Descargar e instalar desde nodejs.org (recomiendo la versión LTS)
        \item Linux:

        \begin{figure}[H]
        \begin{lstlisting}[language=bash]
    sudo apt update
    sudo apt install nodejs npm
        \end{lstlisting}
        \caption{Instalar node.}
        \label{fig:Instalar-node}
        \end{figure}
    \end{itemize}
\end{enumerate}

\subsubsection*{Configurar Front}

\begin{enumerate}
    \item Clonar el repositorio en la carpeta destinada a Front

    \begin{figure}
    \begin{lstlisting}[language=bash]
git clone https://github.com/Rocnarx/Chibchaweb-Hosting-Platform-Frontend.git
    \end{lstlisting}
    \caption{Clonar repo Front.}
    \label{fig:Clonar-repo-Front}
    \end{figure}
     \item Acceder a la carpeta de proyecto e instalar las dependencias

     \begin{figure}[H]
     \begin{lstlisting}[language=bash]
     cd Chibchaweb-Hosting-Platform-Frontend
     npm install
     \end{lstlisting}
     \caption{Instalar dependencias front.}
     \label{fig:req-front}
     \end{figure}
     \item Ejecutar el proyecto en desarrollo

     \begin{figure}[H]
     \begin{lstlisting}[language=bash]
     npm run dev
     \end{lstlisting}
     \caption{Crear base de datos.}
     \label{fig:create-db}
     \end{figure}

     \item Crear archivo .env el cual nos permetira conectar el front con los demás componentes

     \item Configurar archivo .env con el siguiente texto

     \begin{figure}[H]
     \begin{lstlisting}[language=bash]
 VITE_API_URL=https://fastapi-app-production-d0f4.up.railway.app
 VITE_API_KEY=dev-4b872f28-23f9-4a1e-9a60-1cb7f4090abc
 \end{lstlisting}
 \caption{Configurar .env front}
 \label{fig:config-env-front}
    \end{figure}
\end{enumerate}


\newpage
\section{Descripción de los módulos}
%configutación tablas
\pgfplotstableset{
    % Estilo general
    every head row/.style={
        before row=\hline\rowcolor{cafet}\bfseries\color{white},
        after row=\hline,
    },
    every odd row/.style={before row=\rowcolor{gray!10}},
    columns/Método/.style={string type},
    columns/Endpoint/.style={string type, column type=p{7cm}},
    columns/Descripción/.style={string type,column type=p{7cm},},
    col sep=semicolon,
}

\subsection{Backend}
El backend del proyecto fue desarrollado en Python 3.12 bajo FastAPI, utilizando SQLAlchemy ORM para la gestión de la base de datos MySQL, este fue dividido en varios módulos funcionales cómo son:
\subsubsection{Módulo de Gestión de Carrito}
Este módulo permite la gestión de dominios en carritos de compra por parte de la cuenta del cliente, sus endpoint principales se pueden ver en el siguiente cuadro:


\begin{center}
    \begin{table}[H]
    \pgfplotstabletypeset[
    col sep=semicolon,
    header=true,
    ]{desc/Carrito.dat}
    \caption{Módulo de Gestión de Carrito}
    \end{table}
\end{center}


\subsubsection{Módulo de Gestión de Dominios}
Este módulo, cumple la funcionalidad de consultar y registrar dominios, así como verificar disponibilidad de estos y transferirlos por correo electrónico, a continuación se describen los principales endpoin de este:

\begin{center}
    \begin{table}[H]
    \pgfplotstabletypeset[
    col sep=semicolon,
    header=true,
    ]{desc/dominios.dat}
    \caption{Módulo de Gestión de Dominios}
    \end{table}
\end{center}

\subsubsection{Módulo de Cuentas y Autenticación}
Este se centra en la funcionalidad de gestión de todas las cuentas del sistema (Cliente, distribuidor, administrador y empleado), donde se incluye, su registro, su ingreso  y su autenticación segura.

\begin{center}
    \begin{table}[H]
        \pgfplotstabletypeset[
        col sep=semicolon,
        header=true,
        ]{desc/cuentasautenticacion.dat}
        \caption{Módulo de Cuentas y Autenticación}
    \end{table}
\end{center}

\subsubsection{Módulo de Planes y Hosting}
\begin{center}
    \begin{table}[H]
        \pgfplotstabletypeset[
        col sep=semicolon,
        header=true,
        ]{desc/planeshosting.dat}
        \caption{Módulo de Planes y Hosting}
    \end{table}
\end{center}

\subsubsection{Módulo de Soporte y Ticket}
\begin{center}
    \begin{table}[H]
        \pgfplotstabletypeset[
        col sep=semicolon,
        header=true,
        ]{desc/tickets.dat}
        \caption{Módulo de Soporte y Ticket}
    \end{table}
\end{center}

\subsection{Middleware de Automatización (OpenAI)}
El sistema integra un middleware con OpenAI para:
\begin{itemize}
    \item Clasificación automática de solicitudes entrantes (cancelación, reclamo, consulta, agradecimiento, solicitud, otra).
    \item Generación automática de respuesta preliminar y cordial al cliente (vía correo), notificando recepción del ticket y plazo estimado de respuesta (24 a 48 horas).
    \item Gestión de historial de tickets: todos los mensajes se almacenan como archivos JSON para trazabilidad.
\end{itemize}
Este componente utiliza el API de openrouter.ai y las claves definidas en las variables de entorno del backend (OPENAI\_API\_KEY).

\subsection{Base de datos}
La base de datos utilizada en el sistema es de tipo relacional, implementada sobre MySQL, con su gestión y visualización realizada a través de DBeaver. Su diseño responde a la lógica del negocio de un sistema de adquisición de dominios y hosting, integrando funcionalidades de carrito de compras, facturación, métodos de pago y gestión de usuarios.

A continuación, se describe brevemente la estructura general de las entidades más relevantes:\\\\
\textbf{CUENTA:} Representa a los usuarios registrados en el sistema. Almacena la información asociada a la cuenta del cliente.\\\\
\textbf{CARRITO:} Contiene los productos (dominios o paquetes) seleccionados por el usuario para la compra. Se relaciona con las tablas ESTADOCARRITO, CARRITODOMINIO y FACTURA.\\\\
\textbf{DOMINIO:} Define los nombres de dominio disponibles, incluyendo su precio y estado de ocupación.\\\\
\textbf{CARRITODOMINIO}: Tabla intermedia que relaciona los dominios añadidos al carrito.\\\\
\textbf{FACTURA:} Registra los pagos asociados a carritos finalizados, incluyendo el valor total, fecha y método de pago.\\\\
\textbf{PAQUETEHOSING y INFOPAQUETEHOSING:} Juntas describen los servicios de hosting disponibles, sus características técnicas (espacio en disco, certificados SSL, correos, etc.) y periodicidad de cobro.\\\\
\textbf{FACTURAPAQUETE e ITEMPAQUETE:} Vinculan los paquetes de hosting con las facturas emitidas, permitiendo registrar múltiples ítems por factura.\\\\
\textbf{METODOPAGOCUENTA:} Permite al usuario registrar métodos de pago, ya sea tarjeta u otro tipo, asociados a su cuenta.\\\\
\textbf{TARJETA y TIPOTARJETA:} Gestionan los datos de tarjetas de crédito o débito que pueden ser utilizadas como medio de pago.\\\\
\textbf{TIPOMETODOPAGO:} Catálogo que clasifica los diferentes tipos de métodos de pago.\\\\
\textbf{TICKET:} Módulo de soporte para generación y atención de incidencias o solicitudes por parte del cliente.\\\\
\textbf{PAIS:} Entidad referencial para identificar el país del usuario o cuenta.


\newpage
\section{Procedimientos de uso}
A continuación, se detallan los procedimientos básicos para utilizar las funcionalidades del sistema ChibchaWeb, desde la perspectiva técnica. Este apartado no sustituye al manual de usuario final, sino que proporciona una guía operativa para verificar el correcto funcionamiento del sistema y realizar pruebas.

\subsection{Acceso al sistema}
\begin{enumerate}
    \item Abrir el navegador y acceder al dominio principal: \href{https://www.chibchaweb.site/ }{https://www.chibchaweb.site/}
    \item Iniciar sesión utilizando una cuenta registrada:
    \begin{enumerate}
        \item Para pruebas técnicas, se pueden usar cuentas de tipo administrador, distribuidor o empleado. 
        \item Los accesos están restringidos según el idtipocuenta.
    \end{enumerate}
\end{enumerate}

\subsection{Registro de nuevos usuarios}
\begin{enumerate}
    \item Desde el frontend, dirigirse al formulario correspondiente: 
    \begin{enumerate}
        \item Registro de clientes 
        \item Registro de distribuidores
        \item Registro de empleados
    \end{enumerate}
    \item Llenar los campos requeridos: 
    \begin{enumerate}
        \item Nombre, identificación, correo, teléfono, dirección, contraseña, tipo de cuenta, país y plan. 
    \end{enumerate} 
    \item Enviar el formulario. El sistema realiza una solicitud POST a la API (/registrar2). 
    \item Verificar en la base de datos MySQL que los datos se almacenen correctamente en la tabla CUENTA. 
\end{enumerate}

\subsection{Gestión de dominio y facturación}
\begin{enumerate}
    \item Iniciar sesión como usuario autorizado (cliente o distribuidor). 
    \item Navegar a la sección de dominios. 
    \item Realizar una compra o agregar un dominio al carrito. 
    \item Proceder con la facturación. 
    \item Verificar en la base de datos que se actualicen las tablas relacionadas: 
    \begin{enumerate}
        \item DOMINIO, CARRITO, FACTURA, TRANSACCION.
    \end{enumerate} 
\end{enumerate}

\subsection{}{Creación y seguimiento de tickets}
\begin{enumerate}
    \item Desde el frontend, acceder a la sección de soporte/tickets. 
    \item Crear un nuevo ticket describiendo el problema o requerimiento. 
    \item El sistema registra automáticamente la solicitud y, en segundo plano, utiliza la clave de OpenAI desde el backend para generar una respuesta automática inicial (si está habilitado). 
    \item Los tickets pueden ser gestionados desde el panel administrativo. 
\end{enumerate}

\subsection{Verificación de autenticación y permisos}
\begin{itemize}
    \item La autenticación se maneja desde el backend mediante tokens. 
    \item Los permisos se controlan con base en el idtipocuenta asignado. 
    \item Verificar que: 
    \begin{enumerate}
        \item Usuarios con estado idestado = 5 no puedan iniciar sesión. 
        \item Los formularios de login y registro validen correctamente los datos.
    \end{enumerate} 
\end{itemize}






 

\newpage
\section{Glosario y anexos}
\subsubsection*{A}

\subsubsection*{B}

\subsubsection*{C}
\begin{itemize}
\item \textbf{Carrito de compra:} Funcionalidad que permite a los usuarios seleccionar, almacenar y gestionar productos o servicios antes de completar una compra.

\item \textbf{Cliente (cuenta):} Usuario registrado en el sistema que adquiere y gestiona servicios como dominios y hosting.

\item \textbf{Comisión:} Porcentaje o monto fijo asignado a distribuidores por ventas realizadas dentro del sistema.

\item \textbf{Cuenta (CUENTA):} Registro individual de usuario dentro del sistema, ya sea cliente, distribuidor, empleado o administrador.
\end{itemize}

\subsubsection*{D}
\begin{itemize}
    \item \textbf{Distribuidor:} Usuario con permisos especiales para revender servicios de dominio/hosting, gestionando clientes y comisiones.
    \item \textbf{Dominio (DOMINIO):} Nombre único que identifica un sitio en Internet, adquirido y gestionado por usuarios dentro del sistema.
\end{itemize}

\subsubsection*{E}
\begin{itemize}
    \item \textbf{Estado del carrito:} Condición actual del carrito de un usuario (abierto, confirmado, pagado, cancelado, etc.).
    \item \textbf{Estado del ticket:} Seguimiento del progreso de un ticket de soporte (pendiente, en proceso, resuelto, cerrado).

\end{itemize}

\subsubsection*{F}

\subsubsection*{G}
\begin{itemize}
    \item \textbf{Gestión de tickets:} Funcionalidad que permite registrar, asignar y resolver solicitudes o incidencias de usuarios.
    \item \textbf{Gestión de dominios:} Módulo que permite registrar, transferir o administrar dominios vinculados a los clientes.
    \item \textbf{Gestión de cuentas:} Funcionalidad administrativa para crear, editar, suspender o eliminar cuentas del sistema.
\end{itemize}

\subsubsection*{H}
\begin{itemize}
\item \textbf {Hosting:} Servicio de alojamiento web que permite publicar sitios y aplicaciones en Internet, con diferentes planes disponibles.
\end{itemize}

\subsubsection*{I}
\begin{itemize}
    \item \textbf {Identificación de cuenta:} Proceso que permite determinar el tipo, estado y permisos de una cuenta específica.
\end{itemize}

\subsubsection*{J}

\subsubsection*{K}

\subsubsection*{L}

\subsubsection*{M}
\begin{itemize}
    \item \textbf {Método de pago:} Forma en que un cliente puede abonar por servicios (tarjeta, transferencia, etc.).
\end{itemize}

\subsubsection*{N}

\subsubsection*{O}

\subsubsection*{P}
\begin{itemize}
    \item \textbf{Paquete de hosting:} Conjunto de servicios agrupados bajo un plan que incluye almacenamiento, ancho de banda, correos, etc.
    \item \textbf{Plan (de hosting):} Nivel de servicio contratado, con características técnicas y precio definido.
\end{itemize}

\subsubsection*{Q}

\subsubsection*{R}
\begin{itemize}
    \item \textbf{Registro de cuenta:} Proceso por el cual un nuevo usuario se da de alta en el sistema.
    \item \textbf{Roles de usuario:} Clasificación de cuentas según sus privilegios (cliente, distribuidor, etc.).
\end{itemize}

\subsubsection*{S}
\begin{itemize}
    \item \textbf{Soporte técnico:} Área dedicada a resolver consultas, problemas o solicitudes de los usuarios del sistema.
    \item \textbf{Sistema de autenticación:} Mecanismo que valida la identidad de usuarios y otorga acceso según sus permisos.
\end{itemize}

\subsubsection*{T}
\begin{itemize}
\item \textbf{Tarjeta (de pago):} Medio electrónico utilizado por los clientes para abonar productos o servicios.
\item \textbf{Ticket:} Solicitud registrada por un cliente relacionada con soporte técnico, facturación, etc.
\item \textbf{Transferencia de dominio:} Proceso mediante el cual un dominio cambia de registrador o de cuenta dentro del sistema.
\item \textbf{Tipo de cuenta:} Clasificación de una cuenta según su función: cliente, distribuidor, etc.
\item \textbf{Tipo de método de pago:} Categoría asignada al método usado para abonar: tarjeta, efectivo, transferencia.
\end{itemize}

\subsubsection*{U}
\begin{itemize}
\item \textbf{Usuario final:} Persona que utiliza el sistema para contratar servicios y administrar su cuenta.
\end{itemize}

\subsubsection*{V}

\subsubsection*{W}

\subsubsection*{X}

\subsubsection*{Y}

\subsubsection*{Z}

\newpage
\bibliographystyle{plain}
\bibliography{bibliografia/bibliografia}

\end{document}
